% =====================================================================
% 第三讲 (Exposé III) 光滑态射: 延拓性质
% (Morphismes lisses: propriétés de prolongement)
% A. Grothendieck
% 原书页码: pp. 49-73
% =====================================================================
\chapter{光滑态射：延拓性质}

\section{形式光滑同态}

在讲 II 中，我们只限于考虑有限型同态，因此在局部环的局部同态 $A \to B$ 中，只限于 $B$ 同构于 $A$ 上有限型代数的局部化的情形. 这个情形对多种应用 (特别是形式几何 (formal geometry) 或解析几何 (analytic geometry)) 是不够的. 例如，形式幂级数环 $B = A[\![t_1, \dots, t_n]\!]$ 具有 (从形式几何的观点看) $A$ 上光滑代数的性质. 在解析几何中，乘积 $Y \times \mathbf{C}^n$ 的一个点 $(y, z)$ 的局部环亦然，它被看作 $y$ 的局部环上的代数；而且该代数的完备化同构于基环 $\mathcal{O}_x$ 的完备化上的 $n$ 个未定元的形式幂级数代数. 正是这一点导致我们给出下面的定义.

\begin{definition}\label{def:III.1.1}
设 $u \colon A \to B$ 是局部环 (回忆它们是诺特的) 的局部同态. 假设 $k(B)$ 在 $k(A)$ 上有限. 称 $u$ 是形式光滑 (formally smooth) 同态，或称代数 $B$ 在 $A$ 上形式光滑，若存在 $A$ 上的有限局部 $A$-代数 $A'$，在 $A$ 上自由，使得半局部环 $\overline{B} \otimes_{\overline{A}} A' = B'$ 的局部分量作为 $A'$-代数同构于 $A'$ 上的形式幂级数代数\footnote{关于由下面 \ref{thm:III.2.1} 启发的更一般、更概念化的定义，参见 EGA 0$_{\mathrm{IV}}$ 19.3.1.}.
\end{definition}

(用 $\overline{A}$，$\overline{B}$ 记 $A$，$B$ 的完备化.) 由于 $B'$ 在 $\overline{B}$ 上有限且自由，它确是一个半局部环，是完备局部环的直积，其中每个局部环仍是 $\overline{B}$ 上的自由模，故维数与 $\overline{B}$ 相同，从而与 $B$ 相同. 由此得出，\ref{def:III.1.1} 中所考虑的形式幂级数环中变元 $t_i$ 的个数等于 $\dim B - \dim A = \dim \overline{B} - \dim \overline{A}$，特别地与所选择的局部分量无关. 立即看出，它也就是环 $B \otimes k = B/\mathfrak{m}B$ 的维数，其中 $k = A/\mathfrak{m}$ 是 $A$ 的剩余域；称它为 $B$ 相对 $A$ 的\emph{相对维数} (relative dimension).

\begin{remark}\label{rem:III.1.2}
显然，定义 \ref{def:III.1.1} 只依赖于由 $A \to B$ 推出的完备化上的同态 $\overline{A} \to \overline{B}$，这在一定程度上证明了该术语的合理性. 我们在此对定义 \ref{def:I.3.2} b) 和 \ref{def:I.4.1} b) 感到遗憾，它们有误导之嫌，并宁愿在那些定义所考虑的情形说``\emph{形式非分歧}\index[term]{xingshifeifenqi@形式非分歧}'' (formally unramified) 和``\emph{形式平展}\index[term]{xingshipingzhan@形式平展}'' (formally etale)，把``非分歧''和``平展''的术语保留给 $B$ 是 $A$ 上有限型代数的局部化的情形\footnote{或更好地说``本质上非分歧'' (essentially unramified)，相应地``本质上平展'' (essentially etale)，比较 EGA IV 18.6.1.}. 读者立即验证``形式平展''等价于``形式光滑且拟有限''. 最后指出，对剩余扩张 $k(B)/k(A)$ (这里假设有限) 不作任何预先假设，也存在``形式光滑''的合理定义，它尤其涵盖这样的局部同态 $A \to B$：$B$ 在 $A$ 上平坦且 $B/\mathfrak{m}B$ 是 $A/\mathfrak{m} = k$ 的可分扩张 (不必有限型)；例如，Cohen $p$-环在 $p$-adic 整数环上形式光滑. 在这个一般情形，同态的提升性质 (比较 \ref{thm:III.2.1}) 应当成为定义. 对我们所考虑的应用，定义 \ref{def:III.1.1} 所处理的情形已经足够；在下面的``形式光滑''中，我们隐含``剩余扩张有限''这一条件.
\end{remark}

\begin{lemma}\label{lem:III.1.3}
若 $B$ 在 $A$ 上形式光滑，则 $B$ 在 $A$ 上平坦.
\end{lemma}

\begin{proof}
由于平坦性对完备化不变，可以假设 $A$ 和 $B$ 完备. 由于平坦性对基环的平坦 (从而忠实平坦) 局部扩张不变，由定义 \ref{def:III.1.1} 归结到 $B$ 是 $A$ 上的形式幂级数代数的情形. 但此时作为 $A$-模，$B$ 同构于若干同构于 $A$ 的 $A$-模的乘积，故 (基环 $A$ 诺特) 作为平坦 $A$-模的乘积是 $A$-平坦的.
\end{proof}

设我们处于 \ref{def:III.1.1} 的条件下. 由于 $B'$ 的局部分量在 $A'$ 上的剩余扩张平凡，由此得出 $L \otimes_k k'$ 是 $k'$-代数，是 Artin 的，其局部分量的剩余扩张平凡 (其中 $L$，$k$，$k'$ 是 $A$，$B$，$A'$ 的剩余域). 这个有限自由扩张 $A'$ 满足 \ref{def:III.1.1} 所述条件的必要条件也是充分的，这由下面的 \ref{prop:III.1.4} (i) 和 \ref{prop:III.1.5} 立即得出.

\begin{proposition}\label{prop:III.1.4}
设 $A \to B$ 是剩余扩张有限的局部环局部同态，$A'$ 是 $A$ 上的有限局部 $A$-代数，于是 $B' = B \otimes_A A'$ 在 $B$ 上有限，从而也是诺特半局部环. (i) 若 $B$ 在 $A$ 上形式光滑，则 $B$ 在其极大理想处的局部化在 $A'$ 上形式光滑. (ii) 若 $A'$ 在 $A$ 上自由，则逆命题成立.
\end{proposition}

\begin{proof}
立即归结到 $A$，$B$ 完备的情形.
\end{proof}

\begin{enumerate}[(i)]
\item 设 $A''$ 是 $A$ 上的有限自由局部扩张，使得 $B'' = B \otimes_A A''$ 的局部分量是 $A''$ 上的形式幂级数代数. 作标量扩张 $A'' \to A'' \otimes_A A' \to A'''$，其中 $A'''$ 是 $A'' \otimes_A A'$ 的一个局部分量，看到 $B'' \otimes_{A''} A''' = B \otimes_A A'''$ 的局部分量是 $A'''$ 上的形式幂级数代数. 另一方面还有 $B \otimes_A A''' = (B \otimes_A A') \otimes_{A'} A''' = B' \otimes_{A'} A'''$；又由于 $A''$ 在 $A$ 上自由，$A'' \otimes_A A'$ 在 $A'$ 上自由，作为其直因子的 $A'''$ 亦然，这证明 $B'$ 在 $A'$ 上形式光滑.
\item 设 $A''$ 是 $A'$ 上的有限自由局部代数，使得 $B' \otimes_{A'} A'' = B \otimes_A A''$ 的局部分量是 $A''$ 上的形式幂级数代数. 由于 $A'$ 在 $A$ 上自由，$A''$ 也自由，故 $B$ 在 $A$ 上形式光滑.
\end{enumerate}

\begin{proposition}\label{prop:III.1.5}
设 $A \to B$ 是剩余扩张平凡的局部环局部同态. 为使 $B$ 在 $A$ 上形式光滑，必要且充分条件是 $\overline{B}$ 同构于 $\overline{A}$ 上的形式幂级数代数.
\end{proposition}

\begin{proof}
只需证明必要性，且可以假设 $A$ 和 $B$ 完备. 设 $m$ ($n$) 是 $A$ ($B$) 的极大理想，$t_1, \dots, t_n$ 是 $n$ 的元素，定义向量空间
\[
(\mathfrak{n}/\mathfrak{n}^2)/\operatorname{Im}(\mathfrak{m}/\mathfrak{m}^2) = \mathfrak{n}/(\mathfrak{n}^2 + \mathfrak{m}B)
\]
的基. 这些元素因此定义局部 $A$-代数同态
\[
B_1 = A[\![t_1, \dots, t_n]\!] \longrightarrow B,
\]
证明它是同构. 只需证明对 $m$ 的每个幂 $m^q$，模 $\mathfrak{m}^q$ 约化后得到同构 ($B_1$ 和 $B$ 是相应约化环按变动的 $q$ 的投射极限). 由于 $B$ 和 $B_1$ 是 $A$-\emph{平坦模}，故按 $m$-adic 滤链所得的相伴分次环在 $k = A/\mathfrak{m}$ 上由 $\mathrm{gr}(A)$ 张量得到，归结为证明 $B_1/\mathfrak{m}B_1 \to B/\mathfrak{m}B$ 是同构. 考虑到 \ref{lem:III.1.3}，归结到 $A$ 是域 $k$ 的情形. 另一方面，若 $A'$ 是 $A$ 上的有限自由局部代数，使得 $B \otimes_A A'$ 是 $A'$ 上的形式幂级数代数 (N.B. 该代数是局部的，因为 $B$ 在 $A$ 上的剩余扩张平凡)，为证明 $B_1 \to B$ 是同构，只需证明 $B_1 \otimes_A A' \to B \otimes_A A'$ 是同构. 这归结到 $B$ 已经是形式幂级数代数的情形 (本应先作这个约化，再归结到基域是域的情形). 但此时 $B$ 是带代表元域 (field of representatives) $k$ 的正则局部环，而众所周知 (并且由考察 $B_1$ 和 $B$ 上 $\mathfrak{n}_1$-adic 和 $n$-adic 滤链的相伴分次环立即看出) $B_1 \to B$ 是同构，证明完毕.
\end{proof}

\begin{corollary}\label{cor:III.1.6}
若 $B$ 在 $A$ 上形式光滑，则存在 $A$ 上的有限局部 $A$-代数 $A'$，使得 $\overline{B} \otimes_{\overline{A}} \overline{A'} = \overline{(B \otimes_A A')}$ 的局部分量同构于 $\overline{A'}$ 上的形式幂级数代数.
\end{corollary}

\begin{proof}
事实上，若 $L/k$ 是 $B/A$ 的剩余扩张，考虑扩张 $k'/k$，使得 $k'$-代数 $L \otimes_k k'$ 中的剩余扩张平凡. 设 $A'$ 是 $A$ 上的有限自由代数，使得 $A'/\mathfrak{m}A' = k'$ (知道它存在：例如逐步归结到 $k'/k$ 单生成的情形，此时把 $k'$ 在 $k$ 上的生成元的最小多项式的系数提升到 $A$). 它是局部的. 则 $B \otimes_A A'$ 在其极大理想处对 $A'$ 的剩余扩张 $k'$ 的剩余扩张平凡，用 \ref{prop:III.1.5} 完成证明.
\end{proof}

\begin{corollary}\label{cor:III.1.7}
设 $A \to B$ 是局部环的局部同态. 为使 $B$ 在 $A$ 上形式光滑，必要且充分条件是 $B$ 在 $A$ 上平坦且 $B/\mathfrak{m}B$ 在 $k = A/\mathfrak{m}$ 上形式光滑.
\end{corollary}

\begin{proof}
对 $A$ 作适当的有限自由局部扩张 $A'$，用 \ref{prop:III.1.4} (ii)，归结到 $B/A$ 的剩余扩张平凡的情形. 此外，由 \ref{prop:III.1.4} (i) 和 \ref{lem:III.1.3} 已知所述条件是必要的. 对充分性，只需注意 \ref{prop:III.1.5} 的证明在所作的假设下 (假设 $A$ 和 $B$ 完备，这是允许的) 证明 $B$ 是 $A$ 上的形式幂级数代数.
\end{proof}

\begin{remark}\label{rem:III.1.8}
对形式光滑同态，不难发展讲 II 中所研究的光滑态射的一切性质的类似物. 对微分性质，这需要修改 Kähler 微分的通常定义 (参见 I 1)，用完备张量积代替通常张量积. 我们只在此提及这些深渊，因为前面的内容对我们的目的已经足够.
\end{remark}

还要在形式光滑性与讲 II 中发展 (且我们至今尚未使用) 的光滑性概念之间建立联系：

\begin{proposition}\label{prop:III.1.9}
设 $A \to B$ 是局部同态，$B$ 是 $A$ 上有限型代数的局部化. 为使 $B$ 在 $A$ 上光滑，必要且充分条件是它在 $A$ 上形式光滑.
\end{proposition}

\begin{proof}
用 \ref{cor:III.1.7} 和 \ref{thm:II.2.1}，归结到 $A$ 是域的情形.
\end{proof}

用 \ref{prop:III.1.4} (ii) 和 \ref{cor:II.4.13}，$k$ 的适当扩张 $k'$ 把我们归结到 $B/k$ 的剩余扩张平凡的情形. 由 \ref{prop:III.1.5} (相应地：\ref{cor:II.5.2})，$B$ 在 $k$ 上光滑 (相应地：在 $k$ 上形式光滑) 当且仅当 $B$ 是正则局部环 (相应地：其完备化是 $k$ 上的形式幂级数代数). 而这两个条件众所周知是等价的 (剩余扩张平凡时).

\section{形式光滑同态的特征提升性质}

\begin{theorem}\label{thm:III.2.1}
设 $A \to B$ 是定义有限剩余扩张的局部环局部同态. 下列条件等价：
\begin{enumerate}[(i)]
\item $B$ 在 $A$ 上形式光滑；
\item 对每个局部同态 $A \to C$ (其中 $C$ 是完备局部环)，$C$ 的每个含于 $\mathfrak{r}(C)$ 的理想 $J$，以及每个局部 $A$-同态 $B \to C/J$，存在提升它的 (必然局部的) $A$-同态 $B \to C$；
\item 对每个 $A$-代数 $C$ (不必是诺特环)，$C$ 的每个幂零理想 $J$，以及每个连续的 (i.e. 在 $\mathfrak{r}(B)$ 的某个幂上为零的) $A$-同态 $B \to C/J$，存在提升它的 (也必然连续的) $A$-同态 $B \to C$；
\item 同 (ii) 和 (iii) 的陈述，但 $C$ 是在 $A$ 上有限的 Artin 局部环.
\item[(iv bis)] 同 (iv)，但 $J$ 的平方为零.
\end{enumerate}
\end{theorem}

\begin{remarkstar}
对本讲的其余部分，我们只用蕴含 (iv)$\Rightarrow$(i) 或 (iv bis)$\Rightarrow$(i)；直接蕴含 (i)$\Rightarrow$(ii) 将在下一节用另一方法、在 $B$ 是 $A$ 上有限型代数的局部化的情形证明. 回忆在 Cohen 定理的``好''理论\footnote{参见 EGA 0$_{\mathrm{IV}}$ 19.3，19.8.}中，性质 (ii) 或 (iii) 成为形式光滑同态的定义，而 \ref{def:III.1.1} 成为仅在剩余扩张有限的情形成立的一个特征性质. 注意 (ii) 和 (iii) 中任何一个都不比另一个更一般. 可以给出 (等价的) 性质，同时涵盖它们两者：引入线性拓扑化的分离完备环 $C$，$C$ 的闭的拓扑幂零理想，以及连续同态 $A \to C$ (从而把 $C$ 变成拓扑 $A$-代数)；我们把这个修改留给读者.
\end{remarkstar}

 \paragraph*{\ref{thm:III.2.1} 的证明.} 证明 (i)$\Rightarrow$(iii)$\Rightarrow$(ii)，然后 (iv)$\Rightarrow$(i). 由于 (ii)$\Rightarrow$(iv) 平凡，而 (iv) 与 (iv bis) 的等价性通过对满足 $J^n = 0$ 的整数 $n$ 作直接归纳看出，这将完成证明.

(i)$\Rightarrow$(iii). 直接归纳归结到 $J^2 = 0$ 的情形. 由于 $C$ 在 $A$ 上有限，存在 $A$ 的极大理想的幂 $\mathfrak{m}^q$ 零化 $C$. 除以 $\mathfrak{m}^q$，并注意由 \ref{prop:III.1.4} (i)，$B/\mathfrak{m}^q B$ 仍在 $A/\mathfrak{m}^q$ 上形式光滑，可以假设 $A$ Artin. 由 \ref{lem:III.1.3}，$B$ 在 $A$ 上平坦，故 $A$ Artin 时 $B$ 在 $A$ 上自由. 因此存在提升所给同态 $u \colon B \to C/J$ 的 $A$-模同态
\[
w \colon B \longrightarrow C.
\]
令
\[
f(x, y) = w(xy) - w(x)w(y) \quad (x, y \in B),
\]
则 $f(x, y) \in J$，故 $f$ 是 $B \times B$ 到 $J$ 的 $A$-双线性映射. 为使存在提升 $u$ 的代数同态 $v \colon B \to C$，必要且充分条件是存在 $A$-线性映射 $g \colon B \to J$，使得 $v = w + g$ 是代数同态，这写成
\[
\begin{array}{c}
g(1) = 1 - w(1), \\
g(xy) - u(x)g(y) - u(y)g(x) = -f(x, y) \qquad (\text{对 } x, y \in B).
\end{array}
\]
这是 $\mathrm{Hom}_A(B, J)$ 中的线性方程组，右端在 $J$ 中，故它有解当且仅当在 $\mathrm{Hom}_{A}(B, J) \otimes_A A'$ 中、右端在 $J' = J \otimes_A A'$ 中的相应方程组有解，其中 $A'$ 是 $A$ 上的忠实平坦代数. 而设 $A'$ 是 $A$ 上的有限自由局部代数，使得 $B' = B \otimes_A A'$ 是 $A'$ 上的形式幂级数代数 (N.B. 在证明中可以假设 $A$ 和 $B$ 完备，立即看出). 由于 $A'$ 在 $A$ 上有限自由，有
\[
\operatorname{Hom}_{A}(B, J) \otimes_A A' = \operatorname{Hom}_{A'}(B', J'),
\]
并看到在 $\mathrm{Hom}_{A'}(B', J')$ 中得到的方程组，正是确定那些经``修正'' $A'$-模同态 $g' \colon B' \to J'$ 提升由 $u$ 标量扩张得到的 $A'$-模同态 $w' \colon B' \to C'$、从而提升 $u' \colon B' \to C'/J'$ 的 $A'$-代数同态 $B' \to C' = C \otimes_A A'$ 的方程组. (注意 $B$ 作为 $A'$-模生成 $B'$.) 于是归结到证明 $B$ 是 $A$ 上的形式幂级数代数 $B = A[\![t_1, \dots, t_n]\!]$ 时 (iii) 成立. 此时把 $t_i$ 在 $C/J$ 中的像任意提升为 $C$ 的元素 $z_i$. 由于 $z_i$ 模 $J$ 是幂零的 ($u \colon B \to C/J$ 连续)，$z_i$ 亦然 (因为 $J$ 幂零)，故诸 $z_i$ 定义 $B$ 到离散拓扑 $C$ 中的连续拓扑 $A$-代数同态，显然提升 $u$，证毕.

(iii)$\Rightarrow$(ii). 设 $n$ 是 $C$ 的极大理想，对每个整数 $q > 0$，设
\[
C_q = C/\mathfrak{n}^q, \qquad J_q = (J + \mathfrak{n}^q)/\mathfrak{n}^q,
\]
则 $C_q/J_q$ 等同于 $C/J$ 的商代数，另一方面复合同态 $u_q \colon B \to C/J \to C_q/J_q$ 是 $B$ 到离散环 $C_q/J_q$ 的连续同态，且 $J_q$ 是 $C_q$ 中的幂零理想. 于是逐步构造 $A$-同态
\[
v_q \colon B \longrightarrow C_q
\]
使得 (a) $v_q$ 提升 $u_q$，(b) $v_q$ 提升 $v_{q-1}$. 归纳的可行性容易验证，因为 $u_q \colon B \to C/(J + \mathfrak{n}^q)$ 和 $v_{q-1} \colon B \to C/\mathfrak{n}^{q-1}$ 定义同一个同态
\[
B \longrightarrow C/\bigl((J + \mathfrak{n}^q) + \mathfrak{n}^{q-1}\bigr) = C/(J + \mathfrak{n}^{q-1}) = C_{q-1}/J_{q-1},
\]
即 $u_{q-1}$，它们定义一个同态
\[
B \longrightarrow C/J'_q \quad \text{其中} \quad J'_q = (J + \mathfrak{n}^q) \cap \mathfrak{n}^{q-1} \supset \mathfrak{n}^q
\]
(它们都由它约化得到). 归结为把 $B \to C/J'_q$ 提升为 $B$ 到 $C_q$ 的某个含于 $J_q$ 中的理想 $J'_q/\mathfrak{n}^q$ 之商的同态，因此该理想幂零，由假设 (iii) 这是可能的.

做完之后，诸 $v_q$ 定义 $B$ 到 $C_q$ 的投射极限 $C$ 的同态. 由于 $J$ 闭，$J$ 是诸 $J_q$ 的投射极限，故 $v$ 提升 $u$，证毕.

(iv)$\Rightarrow$(i). 首先立即看出，若 (iv) 成立，则对 $A$ 上的有限局部代数 $A'$，它对 $B \otimes_A A'$ 的局部分量在 $A'$ 上仍成立. 取 $A'$ 在 $A$ 上自由且使 $B'$ 在 $A'$ 上的剩余扩张平凡，由 \ref{prop:III.1.4} (ii) 归结到 $B$ 在 $A$ 上的剩余扩张平凡的情形. 此时我们证明一个稍更精确的结果：

\begin{corollary}\label{cor:III.2.2}
在 \ref{thm:III.2.1} 的条件下，再假设 $B$ 在 $A$ 上的剩余扩张平凡. 则 \ref{thm:III.2.1} 的等价条件还等价于下列两个条件 (在 (v) 中假设 $A$ 和 $B$ 完备)：
\begin{enumerate}[(iv ter)]
\item[(iv ter)] 同 (iv)，但 $A$ 上有限的 Artin 局部环 $C$ 限于具有平凡的剩余扩张 (且，如果愿意，理想 $J$ 的平方为零).
\item[(v)] 存在局部 $A$-同态 (其中 $n = \dim \mathfrak{n}/(\mathfrak{n}^2 + \mathfrak{m}B)$)
\[
u \colon B \longrightarrow B_1 = A[\![t_1, \dots, t_n]\!]
\]
诱导同构
\[
\mathfrak{n}/(\mathfrak{n}^2 + \mathfrak{m}B) \stackrel{\sim}{\longrightarrow} \mathfrak{n}_1/(\mathfrak{n}_1^2 + \mathfrak{m}B_1),
\]
其中 $n$ ($n_1$) 是 $B$ ($B_1$) 的极大理想，$m$ 是 $A$ 的极大理想.
\end{enumerate}
\end{corollary}

\paragraph*{证明.} 由于 (iv bis) 显然蕴含 (iv ter) —— 撇开平方为零理想的玩笑不谈 —— 只需证明 (iv ter)$\Rightarrow$(v)$\Rightarrow$(i).

(iv ter)$\Rightarrow$(v). 选择 $\mathfrak{n}/(\mathfrak{n}^2 + \mathfrak{m}B)$ 的基 $a_1, \dots, a_n$，它因此定义局部 $A$-代数同态
\[
B \longrightarrow B_1/(\mathfrak{n}_1^2 + \mathfrak{m}B_1) = k[t_1, \dots, t_n]/(t_1, \dots, t_n)^2,
\]
由 (iv ter)，可以把它逐步提升为 $B$ 到 $B_1/\mathfrak{n}_1^2$，$B_1/\mathfrak{n}_1^3$ 等的 $A$-代数同态，取投射极限得到具有所述性质的同态 $B \to B_1$.

(v)$\Rightarrow$(i). 在交换图中 (两行正合，两端垂直箭头满射)
\[
\begin{tikzcd}[row sep=large, column sep=large]
\mathfrak{n}/(\mathfrak{n}^2 + \mathfrak{m}B) \arrow[r] \arrow[d] & \mathfrak{n}_1/(\mathfrak{n}_1^2 + \mathfrak{m}B_1) \arrow[d] \\
B/\mathfrak{m}B \arrow[r] & B_1/\mathfrak{m}B_1
\end{tikzcd}
\]
中间箭头满射，由此 ($B$ 完备) 得出 $B \to B_1$ 满射. 设 $x_i$ ($1 \leqslant i \leqslant n$) 是 $B$ 中提升 $t_i$ 的元素. 它们因此定义 $A$-代数同态 $B_1 \to B$，它同 $u$ 一样是满射，且按构造与 $u$ 的复合是恒等. 故 $B_1 \to B$ 也单射，从而是同构. 于是得到

\begin{corollary}\label{cor:III.2.3}
在 \ref{cor:III.2.2} (v) 的条件下，$u$ 必然是同构.
\end{corollary}

\begin{proof}
这完成 $B$ 在 $A$ 上形式光滑的证明. 同时还重新得到 \ref{prop:III.1.5} (但这没什么了不起).
\end{proof}

\section{态射到光滑 \texorpdfstring{$S$-概形}{S-概形}中的局部无穷小延拓}

\begin{theorem}\label{thm:III.3.1}
设 $f \colon X \to Y$ 是局部有限型态射. 下列条件等价：
\begin{enumerate}[(i)]
\item $f$ 光滑.
\item 对 $Y$ 上的每个预概形 $Y'$，每个与 $Y'$ 有相同底空间的 $Y'$ 的闭子预概形 $Y'_0$，每个 $Y$-态射 $g_0 \colon Y'_0 \to X$ 以及每个 $z \in Y'_0$，存在 $z$ 在 $Y'$ 中的开邻域 $U$ 以及 $g_0|Y'_0 \cap U$ 到 $Y$-态射 $U \to X$ 的延拓 $g$.
\item[(ii bis)] 对 (ii) 中的 $Y'$，$Y'_0$ 和 $z$，令 $X' = X \times_Y Y'$，$X'_0 = X \times_Y Y'_0$，$X'_0$ 在 $Y'_0$ 上的每个截面延拓为 $X'$ 在 $z$ 的开邻域 $U$ 上的截面.
\item 对每个 $Y$-概形 $Y'$ (为某个 $\mathcal{O}_y$ ($y \in Y$) 上有限的 Artin 局部环的谱)，$Y'$ 的每个非空闭子预概形 $Y'_0$，以及每个 $Y$-态射 $g_0 \colon Y'_0 \to X$，存在延拓 $g_0$ 的 $Y$-态射 $g \colon Y' \to X$.
\item[(iii bis)] 对 (iii) 中的每个 $Y'$，令 $X' = X \times_Y Y'$，$X'_0 = X \times_Y Y'_0$，$X'_0$ 在 $Y'_0$ 上的每个截面延拓为 $X'$ 在 $Y'$ 上的截面.
\end{enumerate}
\end{theorem}

\paragraph*{证明.} (ii) 与 (ii bis) 的等价性，以及 (iii) 与 (iii bis) 的等价性，是平凡的，(ii)$\Rightarrow$(iii) 亦然. 剩下证明 (i)$\Rightarrow$(ii) 和 (iii)$\Rightarrow$(i).

(i)$\Rightarrow$(ii). 设 $x = g_0(z)$. 把 $X$ 换成 $x$ 的适当开邻域，$Y'$ 换成 $g_0$ 在 $x$ 的邻域的原像上的诱导预概形，可以假设 $X$ 在 $Y[t_1, \dots, t_n]$ 上平展. 考虑复合 $Y$-态射 $Y'_0 \to X \to Y[t_1, \dots, t_n]$，它由层 $\mathcal{O}_{Y'_0}$ 的 $n$ 个截面定义，这些截面可以延拓到 $z$ 的邻域成为 $\mathcal{O}_{Y'}$ 的截面，故可以假设该态射已延拓为 $Y$-态射 $Y' \to Y[t_1, \dots, t_n]$. 由 (\ref{cor:I.5.6})，存在唯一提升前者的 $Y$-态射 $g \colon Y' \to X$，它同时延拓 $g_0$，证毕.

(iii)$\Rightarrow$(i). 由于 $f$ 光滑的点集是开的，只需证明它包含每个在纤维中闭的 $x \in X$. 设 $y = f(x)$，则 $\mathcal{O}_x$ 是 $\mathcal{O}_y$ 上的代数，是有限型代数的局部化，剩余扩张有限. 另一方面，假设 (iii) 蕴含：$\mathcal{O}_x$ 到 $A/J$ (其中 $A$ 是 $\mathcal{O}_y$ 上的有限 Artin 局部代数，$J$ 是含于其根式中的理想) 的每个同态，可以提升为 $\mathcal{O}_x$ 到 $A$ 的同态 (考虑到 $\mathrm{Spec}(B)$ ($B$ 局部环) 到 $X$ 的态射由 $\mathcal{O}_x$ ($x \in X$) 到 $B$ 的局部同态一一确定). 由 \ref{thm:III.2.1}，由此得出 $\mathcal{O}_x$ 在 $\mathcal{O}_y$ 上形式光滑，从而由 \ref{prop:III.1.9} 在 $\mathcal{O}_y$ 上光滑.

\begin{corollary}\label{cor:III.3.2}
设 $f \colon X \to Y$ 如 \ref{thm:III.3.1}. 下列条件等价：
\begin{enumerate}[(i)]
\item $f$ 平展；
 \item \ref{thm:III.3.1} 的条件 (ii) 成立，且 $g_0$ 到 $U$ 的延拓 $g$ 唯一；
 \item \ref{thm:III.3.1} 的条件 (iii) 成立，且 $g$ 唯一.
\end{enumerate}
\end{corollary}

\begin{proof}
只需注意，在上述 (i)$\Rightarrow$(ii) 的证明中，只有当 $n = 0$ (已知充分的条件) 时才能有唯一性 (当 $Y'_0$ 在 $z$ 的邻域内不恒等于 $Y'$ 时).
\end{proof}

\begin{corollary}\label{cor:III.3.3}
设 $X$ 是完备局部环 $A$ 上的局部有限型预概形，$y$ 是 $Y = \mathrm{Spec}(A)$ 的闭点，$x$ 是 $f^{-1}(y)$ 中在 $k(y)$ 上有理的点. 若 $X$ 在 $A$ 上于 $x$ 处光滑，则存在 $X$ 在 $Y$ 上的截面 $s$，``经过 $x$''，i.e. 使得 $s(y) = x$.
\end{corollary}

\begin{proof}
特别地，若 $X$ 在 $A$ 上光滑，则自然映射
\[
\Gamma(X/Y) \longrightarrow \Gamma(X \otimes_A k/k)
\]
(从 $X$ 在 $Y$ 上的截面到纤维 $f^{-1}(y) = X \otimes_A k$ 中 $k$ 上有理点的集合) 满射. 这个事实在 $A$ 是离散赋值环且 $X$ 在 $A$ 上紧合 (事实上：射影) 时尤为著名，此时 $X$ 在 $Y$ 上的截面 (i.e.``$X$ 的在 $A$ 中取值的点'') 也等同于有理截面，i.e. $X \otimes_A K = X_K$ (它是 $K$ 上紧合且简单的概形) 中 $K$ 上有理的点 (i.e. $X$ 的 $K$-有理点)，其中 $K$ 是 $A$ 的分式域.
\end{proof}

\section{光滑 \texorpdfstring{$S$-概形}{S-概形}的局部无穷小延拓}

\begin{theorem}\label{thm:III.4.1}
设 $Y$ 是局部诺特预概形，$Y_0$ 是与其有相同底空间的闭子预概形，$X_0$ 是光滑的 $Y_0$-预概形，$x$ 是 $X_0$ 的点. 则存在 $x$ 的开邻域 $U_0$，在 $Y$ 上光滑的预概形 $X$，以及 $Y_0$-同构
\[
h \colon U_0 \stackrel{\sim}{\longrightarrow} X \times_Y Y_0.
\]
此外，若 $(U'_0, X', h')$ 是该问题的另一个解，则``它在 $x$ 的邻域内同构于第一个解''.
\end{theorem}

\begin{proof}
``它''的含义留给读者. 可以注意到，对给定的 $U_0$，该问题的一个解相当于在 $U_0$ 上给出 $f_0^{-1}(\mathcal{O}_Y)$ (其中 $f_0$ 是结构态射 $U_0 \to Y_0$ 之下的连续映射) 的代数层 $\mathcal{B}$，带有与同态 $f^{-1}(\mathcal{O}_Y) \to f^{-1}(\mathcal{O}_{Y_0})$ 相容的环同态 $\mathcal{B} \to \mathcal{O}_{U_0}$，使得
\begin{enumerate}[a)]
\item 该同态诱导同构
\[
\mathcal{B} \otimes_{f^{-1}(\mathcal{O}_Y)} f^{-1}(\mathcal{O}_{Y_0}) \stackrel{\sim}{\longrightarrow} \mathcal{O}_{Y_0};
\]
\item $U_0$ 配备 $\mathcal{B}$ 后成为光滑的 $Y$-预概形.
\end{enumerate}
这样，局部唯一性断言的精确含义变得特别明显.
\end{proof}

\paragraph*{证明.} 已经可以假设 $X_0$ 在 $Y_0[t_1, \dots, t_n] = Y'_0$ 上平展. 而后者可以看作 $Y' = Y[t_1, \dots, t_n]$ 的与它有相同底空间的闭子预概形. 由 (\ref{thm:I.8.3})，存在 $Y'$ 上平展的 $X$ 以及 $Y'_0$-同构 $X \times_{Y'} Y'_0 \stackrel{\sim}{\longrightarrow} X'_0$. 存在性已经解决. 对唯一性，用光滑态射的性质 \ref{thm:III.3.1} (ii)，并考虑下述引理：

\begin{lemma}\label{lem:III.4.2}
设 $Y$ 是预概形，$Y_0$ 是由局部幂零理想层 $\mathcal{J}$ 定义的闭子预概形，$X$ 和 $X'$ 是 $Y$-预概形，$u \colon X \to X'$ 是 $Y$-态射. 假设 $X$ 在 $Y$ 上平坦. 为使 $u$ 是同构，必要且充分条件是 $u_0 \colon X \times_Y Y_0 \to X' \times_Y Y_0$ 是同构.
\end{lemma}

\begin{proof}
证明容易，归结到仿射情形并考察相伴分次环. 还注意到把``同构''换成``闭浸入''得到的类似陈述也成立，且无需平坦性假设.
\end{proof}

\begin{remark}\label{rem:III.4.3}
必须注意， \ref{thm:III.4.1} 中得到的局部延拓 $X$ 不是典范的，换言之，两个解之间的局部同构不唯一，i.e. 一般存在 $X$ 的非平凡 $Y$-自同构，在闭子预概形 $X_0 = X \times_Y Y_0$ 上诱导恒等. 正是因此，在构造光滑预概形的整体无穷小延拓时，应当预期存在上同调性质的\emph{障碍}\index[term]{zhangai@障碍} (obstruction)，它将在下面 (§6) 精确化.
\end{remark}

\section{态射的整体无穷小延拓}

设 $T$ 是拓扑空间，$G$ 是 $X$ 上的群层，$\mathcal{P}$ 是 $T$ 上 $\mathcal{G}$ 在其上作用的集合层 (为确定起见作用在右边). 称 $\mathcal{P}$ 在 $\mathcal{G}$ 下\emph{形式主齐性} (formally principal homogeneous)，若众所周知的集合层同态
\[
\mathcal{G} \times \mathcal{P} \longrightarrow \mathcal{P} \times \mathcal{P}
\]
(由 $\mathcal{G}$ 在 $\mathcal{P}$ 上的作用得出) 是同构. 这等价于说：对每个 $x \in T$，$\mathcal{P}_x$ 是空的，或是在通常群 $\mathcal{G}_x$ 下的主齐性空间；也等价于说：对 $T$ 的每个开集 $U$，$\mathcal{P}(U)$ 是空的，或是在通常群 $\mathcal{G}(U)$ 下的主齐性空间. 称 $\mathcal{P}$ 是 $\mathcal{G}$ 下的主齐性层 (principal homogeneous sheaf)，若它形式主齐性，且诸 $\mathcal{P}_x$ 非空 (换言之，若所有 $\mathcal{P}_x$ 都是 $\mathcal{G}_x$ 下的 (非空的) 主齐性空间)\footnote{看来最好采用 Giraud 论文中引入的更短更形象的术语``$G$ 下的挠子 (torseur under G)''.}. 回忆 $\mathcal{G}$ 下的主齐性层的 (同构) 类的集合等同于上同调集 $H^1(T, \mathcal{G})$，当 $\mathcal{G}$ 交换时它也是 $T$ 以 $\mathcal{G}$ 为系数的通常上同调群. 于是，对每个主齐性 $\mathcal{P}$，有特征类 (characteristic class) $c(\mathcal{P}) \in H^1(T, \mathcal{G})$，其平凡性是 $\mathcal{P}$ 平凡 (i.e. 同构于 $\mathcal{G}$，$\mathcal{G}$ 在其上通过右平移作用) 的必要且充分条件，也就是 $\mathcal{P}$ 有截面的必要且充分条件.

\begin{proposition}\label{prop:III.5.1}
设 $S$ 是预概形，$X$ 和 $Y$ 是 $S$ 上的预概形，$Y_0$ 是由 $Y$ 上平方为零的理想层 $\mathcal{J}$ 定义的 $Y$ 的闭子预概形. 设 $g_0$ 是 $Y_0$ 到 $X$ 的 $S$-态射，$\mathcal{P}(g_0)$ 是 $Y$ 上的层，其在开集 $U$ 上的截面是 $g_0|U \cap Y_0$ 到 $S$-态射 $g \colon U \to X$ 的延拓 $g$. 则 $\mathcal{P}(g_0)$ 是交换群层
\[
\mathcal{G} = \mathbf{Hom}_{\mathcal{O}_{Y_0}}\bigl(g_0^{*}(\varOmega_{X/S}^{1}), \mathcal{J}\bigr)
\]
下的 (自然地) 形式主齐性层.
\end{proposition}

\begin{proof}
令 $\mathcal{P} = \mathcal{P}(g_0)$. 必须对 $Y$ 的每个开集 $U$ 定义映射
\[
\mathcal{P}(U) \times \mathcal{G}(U) \longrightarrow \mathcal{P}(U),
\]
使得 (a) 对固定的 $g \in \mathcal{P}(U)$，映射 $s \mapsto gs$ 从 $\mathcal{G}(U)$ 到 $\mathcal{P}(U)$ 是双射；(b) $\mathcal{P}(U)$ 成为以 $\mathcal{G}(U)$ 为算子群的集合；(c) 前面的映射与 $V \subset U$ 的限制算子相容. (c) 的验证平凡，为简化可以假设 $U = Y$. (b) 的验证 (如所愿是局部性质的) 留给读者，因此对给定的 $g \in \mathcal{P}(Y)$，只定义从 $\mathcal{G}(Y)$ 到 $\mathcal{P}(Y)$ 的自然双射. 于是假设已给出 $S$-态射 $g \colon Y \to X$，求典范双射
\begin{equation}
\mathrm{Hom}_{\mathcal{O}_{Y_0}}\bigl(g_0^{*}(\varOmega_{X/S}^{1}), \mathcal{J}\bigr) \stackrel{\sim}{\longrightarrow} \mathcal{P}(Y), \tag{($*$)}
\end{equation}
其中 $\mathcal{P}(Y)$ 是与 $g$ 诱导同一个态射 $g_0 \colon Y_0 \to X$ 的 $Y$ 到 $X$ 的 $S$-态射 $g'$ 的集合. 这样的 $g'$ 的给出等价于 $S$-态射 $h \colon Y \to X \times_S X$ 的给出，满足 $\mathrm{pr}_1 \circ h = g$，且 $h \circ i = (g_0, g_0)$，其中 $\mathrm{pr}_1 \colon X \times_S X \to X$ 是第一投影，$i \colon Y_0 \to Y$ 是典范浸入，$(g_0, g_0) \colon Y_0 \to X \times_S X$ 是分量为 $g_0$，$g_0$ 的态射 $\Delta_{X/S} g_0$：
\[
\begin{tikzcd}[row sep=large, column sep=large]
X \times_S X \arrow[d,"\mathrm{pr}_1"] & Y_0 \arrow[l,"{h_0 = (g_0, g_0)}"] \\
X & Y \arrow[l,"h"]
\end{tikzcd}
\]
由于 $h_0$ 通过对角浸入 $\Delta_{X/S}$ 分解，且 $Y$ 在 $Y_0$ 的一阶无穷小邻域中 (i.e. $\mathcal{J}^2 = 0$)，所求的 $h$ 必然 (唯一地) 通过对角线的一阶无穷小邻域分解，后者 (作为 $X$-预概形，借助 $\mathrm{pr}_1$) 等同于代数层 $\mathcal{O}_X + \varOmega_{X/S}^{1}$ 的谱 $X'$ (其中第二项看作平方为零的理想)，对角态射 $X \to X'$ 对应于该代数层的典范增广. 令 $Y' = X' \times_X Y$，$Y'_0 = Y' \times_Y Y_0 = X' \times_X Y_0$，于是所求的 $h$ 与 $Y'$ 在 $Y$ 上的、延拓 $Y'_0$ 在 $Y_0$ 上的给定截面 $u_0$ 的截面 $u$ 一一对应. 此外可以把 $Y'$ 等同于 $Y$ 上的代数层 $\mathcal{A} = g^{*}(\mathcal{O}_X + \varOmega_{X/S}^{1}) = \mathcal{O}_Y + g^{*}(\varOmega_{X/S}^{1})$ 的谱，把 $Y'_0$ 等同于代数层 $\mathcal{A}_0 = \mathcal{A} \otimes_{\mathcal{O}_Y} \mathcal{O}_{Y_0} = \mathcal{O}_{Y_0} + g_0^{*}(\varOmega_{X/S}^{1})$ 的谱，则截面 $u_0$ 是由 $A_0$ 到 $\mathcal{O}_{Y_0}$ 的典范增广定义的截面. 故 $\mathcal{P}(Y)$ 等同于诱导典范增广 $\mathcal{A}_0 \to \mathcal{O}_{Y_0}$ 的代数同态 $\mathcal{A} \to \mathcal{O}_Y$ 的集合. 而代数同态 $\mathcal{A} \to \mathcal{O}_Y$ 与模层同态 $\mathcal{M} \to \mathcal{A}$ (为简化令 $\mathcal{M} = g^{*}(\varOmega_{X/S}^{1})$) 一一对应，且所关注的是那些诱导零同态 $\mathcal{M}_0 \to \mathcal{O}_{Y_0}$ (其中 $\mathcal{M}_0 = \mathcal{M} \otimes_{\mathcal{O}_Y} \mathcal{O}_{Y_0}$)，i.e. 把 $\mathcal{M}$ 映入增广理想 $\mathcal{J}$ 的同态. 于是得到集合
\[
\mathrm{Hom}_{\mathcal{O}_Y}(\mathcal{M}, \mathcal{J}) = \mathrm{Hom}_{\mathcal{O}_{Y_0}}(\mathcal{M}_0, \mathcal{J})
\]
(因为 $\mathcal{J}$ 被 $\mathcal{J}$ 零化). 这就是所求的典范双射 ($*$).
\end{proof}

\begin{corollary}\label{cor:III.5.2}
若 $X$ 在 $S$ 上光滑 (至少在 $g_0(Y_0)$ 的点处)，则 $\mathcal{P}$ 甚至是交换群层 $G$ 下的主齐性层，此时 $G$ 也可以写成
\[
\mathcal{G} = g_0^{*}(\mathfrak{g}_{X/S}) \otimes_{\mathcal{O}_{Y_0}} \mathcal{J},
\]
其中 $\mathfrak{g}_{X/S}$ 是 $\varOmega_{X/S}^{1}$ 的对偶层，i.e. $X$ 相对 $S$ 的切层 (tangent sheaf) (或导子层). (最后一个公式来自 $\varOmega_{X/S}^{1}$ 此时有限自由这一事实.)
\end{corollary}

\begin{proof}
特别地，与这个主齐性层相对应的是 $H^1(Y_0, \mathcal{G})$ 中的一个上同调类，其消没是存在延拓 $g_0$ 的 $S$-态射 $g$ 的必要且充分条件. 若存在这样的延拓，则一切可能延拓的集合是群 $H^0(Y_0, \mathcal{G})$ 下的齐性空间.
\end{proof}

在应用形式几何的方法时，情形通常如下：给出两个 $S$-预概形 $X$ 和 $Y$，$S$ 上的凝聚理想层 $\mathcal{I}$；用 $S_n$ 记由 $\mathcal{I}^{n+1}$ 定义的 $S$ 的闭子预概形，并令
\[
X_n = X \times_S S_n, \qquad Y_n = Y \times_S S_n.
\]
假设给出 $S_n$-态射
\[
g_n \colon Y_n \longrightarrow X_n
\]
(或等价地，$S$-态射 $Y_n \to X$，或 $S_{n+1}$-态射 $Y_n \to X_{n+1}$，因为这样的态射必然诱导 $Y_n \to X_n$)，并设法把它延拓为 $S_{n+1}$-态射
\[
g_{n+1} \colon Y_{n+1} \longrightarrow X_{n+1}.
\]
(若可以无限继续，就得到 $Y$ 和 $X$ 关于理想 $\mathcal{I}\mathcal{O}_Y$ 和 $\mathcal{I}\mathcal{O}_X$ 的形式完备化预概形 $\widehat{Y} \to \widehat{X}$ 的态射.) 可以应用 \ref{prop:III.5.1}，把 $(S, X, Y, Y_0, g_0)$ 换成 $(S_{n+1}, X_{n+1}, Y_{n+1}, Y_n, g_n)$，此时群层 $G$ 成为 $g_n^{*}\bigl(\varOmega_{X_{n+1}/S_{n+1}}^{1}\bigr)$ 到 $\mathcal{J} = \mathcal{I}^{n+1}\mathcal{O}_Y/\mathcal{I}^{n+2}\mathcal{O}_Y$ 的模层同态层. 由于 $\mathcal{J}$ 被 $\mathcal{I}\mathcal{O}_Y$ 零化，可以把 $g_n^{*}\bigl(\varOmega_{X_{n+1}/S_{n+1}}^{1}\bigr)$ 换成它在 $Y_0$ 上诱导的层，即 $h_0^{*}\bigl(\varOmega_{X/S}^{1}\bigr)$，其中 $h_0$ 是复合 $Y_0 \to Y_n \to X_{n+1}$，或复合 $Y_0 \to X_0 \to X_{n+1}$，其中 $g_0 \colon Y_0 \to X_0$ 由 $g_n$ 诱导. 由于 $\varOmega_{X_{n+1}/S_{n+1}}^{1}$ 在 $X_0 = X_{n+1} \times_{S_{n+1}} S_0$ 上的原像为 $\varOmega_{X_0/S_0}^{1}$，立即看到
\[
\mathcal{G} = \mathbf{Hom}_{\mathcal{O}_{Y_0}}\bigl(g_0^{*}(\varOmega_{X_0/S_0}^{1}), \mathcal{I}^{n+1}\mathcal{O}_Y/\mathcal{I}^{n+2}\mathcal{O}_Y\bigr).
\]
于是得到：

\begin{corollary}\label{cor:III.5.3}
设 $S$，$X$，$Y$，$\mathcal{I}$，$g_n$ 如上，$\mathcal{P}(g_n)$ 是 $Y$ 上的层，其在开集 $U$ 上的截面是把 $g_n$ 延拓为 $S_{n+1}$-态射 $Y_{n+1} \to X_{n+1}$ 的 $g_{n+1}$. 则 $\mathcal{P}(g_n)$ 是群层
\[
\mathcal{G} = \mathbf{Hom}_{\mathcal{O}_{Y_0}}\bigl(g_0^{*}(\varOmega_{X_0/S_0}^{1}), \mathrm{gr}_{\mathcal{I}\mathcal{O}_Y}^{n+1}(\mathcal{O}_Y)\bigr)
\]
下的形式主齐性层.
\end{corollary}

特别地：

\begin{corollary}\label{cor:III.5.4}
若此外 $X$ 在 $S$ 上光滑 (至少在 $g_0(Y_0)$ 的点处)，则 $\mathcal{P}(g_n)$ 甚至是主齐性层. 特别地，它定义 $H^1(Y_0, \mathcal{G})$ 中的一个障碍类，其消没是存在 $g_n$ 的整体延拓 $g_{n+1}$ 的必要且充分条件. 若存在这样的延拓，则一切整体延拓的集合是 $H^0(Y_0, \mathcal{G})$ 下的主齐性空间. 最后，在所考虑的情形，层 $\mathcal{G}$ 也可以写成
\[
\mathcal{G} = g_0^{*}(\mathfrak{g}_{X_0/S_0}) \otimes_{\mathcal{O}_{Y_0}} \mathrm{gr}_{\mathcal{I}\mathcal{O}_Y}^{n+1}(\mathcal{O}_Y).
\]
\end{corollary}

\begin{proof}
逐步进行，看到若所有 $H^1(Y_0, \mathcal{G}_n)$ 都消没 (其中 $\mathcal{G}_n = g_0^{*}\bigl(\mathfrak{g}_{X_0/S_0}\bigr) \otimes \mathrm{gr}_{\mathcal{I}\mathcal{O}_Y}^{n}(\mathcal{O}_Y)$)，则从任意的 $g_k$ 出发，可以把它依次延拓为 $g_{k+1}$，$\dots$. 特别地，若 $\mathcal{I}$ 幂零，就能找到 $g_k$ 到 $Y$ 的延拓 $g$. $H^1$ 消没的条件在 $Y_0$ 仿射时尤其成立. 于是得到：
\end{proof}

\begin{corollary}\label{cor:III.5.5}
在定理 \ref{thm:III.3.1} 的陈述中，得到一个与其余条件等价的必要且充分条件：假设 (ii) (或 (ii bis)) 中出现的 $Y'$ 仿射，并要求 $g_0$ 存在到整个 $Y'$ 的整体延拓 $g$.
\end{corollary}

注意 \ref{thm:III.3.1} 的证明不能直接给我们这个结果.

一个重要情形是 $Y$ 在 $S$ 上平坦，此时有
\[
\mathrm{gr}^{n}(\mathcal{O}_Y) = \mathrm{gr}^{n}(\mathcal{O}_S) \otimes_{\mathcal{O}_{S_0}} \mathcal{O}_{Y_0},
\]
且当诸 $\mathrm{gr}^{n}(\mathcal{O}_S)$ 在 $S$ 上局部自由时，得到
\[
\mathcal{G}_n = \mathbf{Hom}_{\mathcal{O}_{Y_0}}\bigl(g_0^{*}(\varOmega_{X_0/S_0}^{1}), \mathcal{O}_{Y_0}\bigr) \otimes_{\mathcal{O}_{S_0}} \mathrm{gr}^{n}(\mathcal{O}_S),
\]
或，若 $\varOmega_{X_0/S_0}^{1}$ 也局部自由 (例如 $X$ 在 $S$ 上光滑)
\[
\mathcal{G}_n = g_0^{*}\bigl(\mathfrak{g}_{X_0/S_0}\bigr) \otimes_{\mathcal{O}_{S_0}} \mathrm{gr}^{n}(\mathcal{O}_S).
\]
例如若 $S$ 仿射，环为 $A$，$\mathcal{I}$ 由 $A$ 的理想 $I$ 定义，得到
\[
H^{i}(Y_0, \mathcal{G}_n) = H^{i}(Y_0, \mathcal{G}_0) \otimes_A \mathrm{gr}_{I}^{n}(A)
\]
对每个 $i$ 成立 (事实上，问题在 $S_0$ 上是局部的，归结为与自由模层张量的情形). 此时，$H^1(Y_0, \mathcal{G}_0)$ 的消没蕴含 $g_n$ 的一切逐次延拓障碍都为零. 于是得到：

\begin{corollary}\label{cor:III.5.6}
设 $(S, X, Y, \mathcal{I}, g_n)$ 如上，再假设 $X$ 在 $S$ 上光滑，$Y$ 在 $S$ 上平坦，最后 $S$ 仿射，且诸 $\mathrm{gr}^{n}(\mathcal{O}_S) = \mathcal{I}^n/\mathcal{I}^{n+1}$ 局部自由. 则构造 $g_{n+1}$ 的障碍位于 $H^1(Y_0, \mathcal{G}_0) \otimes_A \mathrm{gr}_{I}^{n+1}(A)$ (其中 $A$ 是 $S$ 的环，$I$ 是定义 $\mathcal{I}$ 的 $A$ 的理想)，只要令
\[
\mathcal{G}_0 = g_0^{*}(\mathfrak{g}_{X_0/S_0}).
\]
若 $H^1(Y_0, \mathcal{G}_0) = 0$，则 $g_n$ 可以延拓为 $S$-态射 $\widehat{g} \colon \widehat{Y} \to \widehat{X}$.
\end{corollary}

当然，如果出发点不是通常的 $S$-预概形 $X$ 和 $Y$，而是 $\widehat{S}$-adic 形式预概形 $X$ 和 $Y$，这个结果原样仍然成立. 例如，它可以用来证明某些在完备局部环上紧合的形式概形 (例如) 实际上是代数的. 事实上，像引理 \ref{lem:III.4.2} 那样进行，得到：

\begin{corollary}\label{cor:III.5.7}
在 \ref{cor:III.5.6} 的条件下，若 $g_0$ 是同构，则 $\widehat{g}$ 亦然.
\end{corollary}

(N.B. 对闭浸入也有同样的结果.)

于是得到：

\begin{proposition}\label{prop:III.5.8}
设 $A$ 是极大理想为 $m$、剩余域为 $k$ 的完备局部环，$X$ 和 $Y$ 是 $A$ 上的两个 $m$-adic 形式预概形，在 $A$ 上平坦 (i.e. 对每个 $n$，$X_n$ 和 $Y_n$ 在 $A_n = A/\mathfrak{m}^{n+1}$ 上平坦)，假设 $X_0 = \mathfrak{X} \otimes_A k$ 在 $k$ 上光滑且 $H^1(X_0, \mathfrak{g}_{X_0/k}) = 0$. 则 $Y_0$ 到 $X_0$ 的每个 $k$-同构延拓为 $Y$ 到 $X$ 的 $A$-同构；若还有 $H^0(X_0, \mathfrak{g}_{X_0/k}) = 0$，则这个延拓唯一.
\end{proposition}

\begin{proof}
这特别给出 $A$ 上光滑形式预概形按给定预概形 $X_0$ 约化的唯一性结果 (在 $H^1(X_0, \mathfrak{g}_{X_0/k}) = 0$ 的条件下). 此外，若 $X$ 和 $Y$ 来自 $A$ 上紧合的普通概形 $X$ 和 $Y$，则由形式几何中层存在性定理 (参见 Bourbaki 讨论班报告，N$^{\circ}$182\footnote{证明参见 EGA III 5.4.1.})，知道 $A$-同构 $Y \xrightarrow{\sim} X$ 与形式完备化的 $A$-同构之间有一一对应，因此
\end{proof}

\begin{corollary}\label{cor:III.5.9}
前面的陈述 \ref{prop:III.5.8} 在把 $X$ 和 $Y$ 换成 $A$ 上紧合的普通 $A$-概形 $X$ 和 $Y$ 时仍然成立.
\end{corollary}

最后，当 $X$ 是 $A$ 上紧合的形式概形，且 $Y$ 形如 $\widehat{Y}$ (其中 $Y$ 是 $A$ 上紧合的普通概形) 时，命题 \ref{prop:III.5.8} 给出 $X$ 与 $\widehat{Y}$ 同构的充分条件，因而形式概形 $X$ 事实上``代数化'' (i.e. 同构于 $\widehat{X}$，其中 $X$ 是 $A$ 上紧合的普通概形，后者将被典范地确定). 例如当 $X_0 = \mathbf{P}_k^r$ 时 (或更一般地，当 $X_0$ 是 Severi-Brauer 概形，i.e. 在 $k$ 的代数闭包上同构于标准射影空间时) 就如此：$A$ 上以 $\mathbf{P}_k^r$ 为纤维的每个紧合平坦形式概形都可代数化，更精确地说同构于 $\mathbf{P}_A^r$ 的 $m$-adic 形式完备化. 特别地 (借助``存在性定理'')，$A$ 上以 $\mathbf{P}_k^r$ 为纤维的每个紧合普通概形同构于 $\mathbf{P}_A^r$ ($A$ 是完备局部环). 用下降理论可以证明，若 $A$ 不完备，则 $X$ 在作基的有限平展扩张 $A \to A'$ 后同构于 $\mathbf{P}^r$ (在这种形式下，结果对 Severi-Brauer 概形纤维仍然成立).

\section{光滑 \texorpdfstring{$S$-概形}{S-概形}的整体无穷小延拓}

在定理 \ref{thm:III.4.1} 的条件下，问题是要知道是否存在 $Y$ 上光滑的预概形 $X$，使得 $X \times_Y Y_0$ 是 $Y_0$-同构于 $X_0$ 的，已知这样的概形``在 $X_0$ 上局部存在''. 沿用逐步构造的方法，把 $Y$ 换成字母 $S$，假设给出由理想层 $\mathcal{I}$ (不再需要假设局部幂零) 定义的 $S$ 的闭子预概形 $S_0$，引入由 $\mathcal{I}^{n+1}$ 定义的 $S$ 的闭子预概形 $S_n$，并假设给出了在 $S_n$ 上光滑的子预概形 $X_n$. 目标是找到``按 $X_n$ 约化''的 $S_{n+1}$-预概形 $X_{n+1}$，i.e. 配备同构
\[
X_{n+1} \times_{S_{n+1}} S_n \stackrel{\sim}{\longleftarrow} X_n,
\]
且在 $S_{n+1}$ 上光滑 (或由 \ref{thm:II.2.1}，等价地在 $S_{n+1}$ 上平坦). 如 §4 所指出的，这样的给出相当于在 $f^{-1}(\mathcal{O}_{S_{n+1}})$ (其中 $f$ 是结构态射 $X_n \to S_n$ 之下的连续映射) 上给出代数层 $\mathcal{B}$，配备增广 $\mathcal{B} \to \mathcal{O}_{X_n}$，与增广 $f^{-1}(\mathcal{O}_{S_{n+1}}) \to f^{-1}(\mathcal{O}_{S_n})$ 相容，并满足 (我们不重写，只指出它们是 $X_n$ 的底空间上的局部性质) 两个条件 (a) 和 (b). 由 \ref{thm:III.4.1} 知道一个解局部存在. 它还是同构 (不唯一) 意义下唯一的，至少在局部如此. 先精确化这一点：

\begin{proposition}\label{prop:III.6.1}
设 $X_{n+1}$ 是在 $S_{n+1}$ 上按 $S_n$ 上的 $X_n$ 约化的. 则 ($X_n$ 的底空间上，或 $X_0$ 上) 在 $X_n$ 上诱导恒等的 $X_{n+1}$ 的 $S_{n+1}$-自同构层，典范同构于
\[
\mathcal{G} = \mathfrak{g}_{X_0/S_0} \otimes_{\mathcal{O}_{S_0}} \mathrm{gr}_{\mathcal{I}}^{n+1}(\mathcal{O}_S)
\]
(作为群层).
\end{proposition}

\begin{proof}
事实上，由 \ref{cor:III.5.4} 和 \ref{lem:III.4.2}，这个层是 $\mathcal{G}$ 下的主齐性层. 由于它带有特惠截面 (即 $X_{n+1}$ 的恒等自同构)，它因此作为集合层等同于 $G$. 必须验证这个等同与群结构相容. 这是容易的，而且它是 \ref{prop:III.5.1} 和 \ref{cor:III.5.3} 中主丛结构关于态射复合的相容性的更一般结果的特殊情形 (我们在此不陈述该结果，但它应当载入大全 (hyperplodoque)).
\end{proof}

特别地，$X_0$ 上 $X_{n+1}$ 的自同构芽层 (带所述结构) 是交换的. 由此得出，若 $X'_{n+1}$ 是该问题的另一个解，在 $X_0$ 的开集 $U$ 之上同构于 $X_{n+1}$，则由同构 $X_{n+1}|U \stackrel{\sim}{\longrightarrow} X'_{n+1}|U$ 通过结构转移得到的 $\mathbf{Aut}(X_{n+1})|U$ 到 $\mathbf{Aut}(X'_{n+1})|U$ 的同构不依赖于后者的选择. (它无非是 $\mathcal{G}$ 的恒等同构，当我们借助 \ref{prop:III.6.1} 把两个自同构层都等同于 $G$ 时.)

由 \ref{prop:III.6.1} 推出：

\begin{corollary}\label{cor:III.6.2}
设 $X_{n+1}$，$X'_{n+1}$ 在 $S_{n+1}$ 上光滑且``按 $X_n$ 约化''. 则在 $X_n$ 上诱导恒等的 $X_{n+1}$ 到 $X'_{n+1}$ 的 $S_{n+1}$-同构层 (在 $X_0$ 的底空间上) 自然地是 $\mathcal{G}$ 下的主齐性层.
\end{corollary}

\begin{proof}
这确实表示 $X_{n+1}$ 与 $X'_{n+1}$ 局部同构，且第一个的自同构芽层为 $\mathcal{G}$.
\end{proof}

现在注意，由 \ref{thm:III.4.1}，总可以找到 $X_n$ 的开覆盖 $(U_i)$ (可以假设它们仿射)，以及每个 $i$ 在 $S_{n+1}$ 上按 $U_i$ 约化的光滑概形 $X^i$. 为简化假设 $X_n$ 分离，故 $U_{ij} = U_i \cap U_j$ 仍是 $X_n$ 的仿射开集. 由于仿射开集的以拟凝聚层 $\mathcal{G}$ 为系数的 $H^1$ 为零，由推论 \ref{cor:III.6.2} 得出 $X^i|U_{ij}$ 同构于 $X^j|U_{ij}$；设
\[
f_{ji} \colon X^i|U_{ij} \xrightarrow{\sim} X^j|U_{ij}
\]
是这样一个同构. 它由 $\mathcal{G}$ 在 $U_{ij}$ 上的一个截面差一个差确定. 对每个三元指标组令：
\[
f_{ji}^{(k)} = f_{ji}|U_{ijk}, \quad \text{其中 } U_{ijk} = U_i \cap U_j \cap U_k.
\]
若成立
\begin{equation}
f_{kj}^{(i)} f_{ji}^{(k)} = f_{ki}^{(j)}, \tag{1}
\label{eq:III.1}
\end{equation}
则诸 $X^i$ 可以通过 $f_{ji}$``粘合''，因此它们定义所求问题的解 $X = X_{n+1}$. 更一般地，若可以把诸 $f_{ji}$ 修改为
\begin{equation}
f'_{ji} = f_{ji} g_{ji} \quad (g_{ji} \in \Gamma(U_{ij}, \mathcal{G})) \tag{2}
\label{eq:III.2}
\end{equation}
使得诸 $f'_{ji}$ 满足上面的传递条件，则这样的解存在. 解存在的这个充分条件也是必要的，回忆：这样的解 $X$ 必须在每个 $U_i$ 上同构于 $X^i$，这使我们能够选择同构
\[
f_i \colon X|U_i \stackrel{\sim}{\longrightarrow} X^i,
\]
并定义满足粘合条件的
\[
f'_{ji} = (f_j|U_{ij})(f_i|U_{ij})^{-1} \colon X^i|U_{ij} \xrightarrow{\sim} X^j|U_{ij}.
\]

令
\begin{equation}
f_{ijk} = (f_{ki}^{(j)})^{-1} f_{kj}^{(i)} f_{ji}^{(k)}, \tag{3}
\label{eq:III.3}
\end{equation}
这是 $X^i|U_{ijk}$ 的自同构，借助 \ref{prop:III.6.1} 把它等同于 $\mathcal{G}$ 的截面. 通过留给读者的小小形式计算，看到它是开覆盖 $\mathcal{U} = (U_i)$ 以 $\mathcal{G}$ 为系数的 2-上闭链 $f$. 同样的计算表明，在 (2) 的条件下，对 $f'_{ij}$ 的粘合条件 (1) 等价于公式
\begin{equation}
f = dg, \tag{4}
\label{eq:III.4}
\end{equation}
其中 $g = (g_{ij})$ 看作 $\mathcal{U}$ 以 $\mathcal{G}$ 为系数的 1-上链. 故问题有解的必要且充分条件是上闭链 (3) 定义的 $H^2(\mathcal{U}, \mathcal{G})$ 中的上同调类为零. 还注意到，由于 $\mathcal{U} = (U_i)$ 是概形 $X_0$ 的仿射覆盖，$H^2(\mathcal{U}, \mathcal{G})$ 等同于 $H^2(X_0, \mathcal{G})$. 显然，在 $H^2(X_0, \mathcal{G})$ 中这样得到的上同调类不依赖于所考虑的仿射覆盖. 称它为把 $X_n$ 延拓为 $S_{n+1}$ 上光滑概形 $X_{n+1}$ 的障碍类.

假设这个障碍为零. 则上面概述的推理表明，每个解 $X = X_{n+1}$ 同构于由同构 $f'_{ji}$ 粘合得到的解，可以假设 $f'_{ji}$ 形如 (2)，粘合条件正是 (3). 故允许的 $g$ 的集合是以 $\mathcal{U}$ 以 $\mathcal{G}$ 为系数的 1-上闭链群 $Z^1(\mathcal{U}, \mathcal{G})$ 为主齐性空间. 此外立即看到，两个上链 $g$ 和 $g'$ (满足 $dg = dg' = f$) 定义同构的解，当且仅当上闭链 $g - g'$ 形如 $dh$，其中 $h = (h_i) \in C^0(\mathcal{U}, \mathcal{G})$. 于是得到：

\begin{theorem}\label{thm:III.6.3}
设 $(S, \mathcal{I}, X_n)$ 如上，假设 $X_n$ 分离\footnote{这个条件实际上不必要，可以避免上面的上闭链计算. 参见 J. Giraud 的书《Cohomologie Non Abélienne》(将出版于 Springer Verlag 1971). 比较注记 \ref{rem:III.6.4}.}. 则可以典范地定义 $H^2(X_0, \mathcal{G})$ (其中 $\mathcal{G}$ 在 \ref{prop:III.6.1} 中定义) 中的障碍类，其消没是存在按 $X_n$ 约化的、在 $S_{n+1}$ 上光滑的概形 $X_{n+1}$ 的必要且充分条件. 若这个障碍为零，则按 $X_n$ 约化的 $S_{n+1}$-预概形 $X_{n+1}$ 的 (在 $X_n$ 上诱导恒等的) 同构类集合，自然地是 $H^1(X_0, \mathcal{G})$ 下的主齐性空间.
\end{theorem}

\begin{remark}\label{rem:III.6.4}
从 \ref{prop:III.6.1} 出发，这里所做的推理纯属形式，并且可以有利地搬到局部范畴甚至一般纤维范畴的框架中. 在一般语境下这样得到的``整体''对象存在性的障碍类 (对其可以局部找到一个对象，且任意两个对象总是``局部同构''，每个对象的自同构群交换) 包含作为特殊情形：不必交换的群层正合序列中的``第二个边缘同态'' (例如 Grothendieck 在 Kansas 或 Tôhoku 中所研究的). 这里所做的笨拙上闭链计算应被看作权宜之计，因为尚不存在令人满意的参考文献.
\end{remark}

\refstepcounter{theorem}\noindent\textbf{\thesection.\arabic{theorem}.}\quad 注意，在 \ref{thm:III.6.3} 中，一般没有 $H^1(X_0, \mathcal{G})$ 之下所考虑主齐性空间中的特惠元素. 这尤其表现为：得到 (在 $S$ 上局部化) $S_0$ 上结构群为 $R^1 f_{*}(\mathcal{G})$ 的主齐性层，它不必平凡，i.e. 定义 $H^1(S_0, R^1 f_{*}(\mathcal{G}))$ 中不必为零的上同调类. (当假设类 $d \in H^2(X_0, \mathcal{G})$ 不为零、但在``$S$ 之上局部为零''，i.e. 定义 $R^2 f_{*}(\mathcal{G})$ 的零截面，i.e. $H^0(S_0, R^2 f_{*}(\mathcal{G}))$ 中的零元素时，正是如此.)

\refstepcounter{theorem}\noindent\textbf{\thesection.\arabic{theorem}.}\quad 目前对本节所引入上同调类的一般代数机制、以及它与前一节的关系几乎一无所知，即使在最简单的特殊情形 (例如 Artin 环上的阿贝尔概形的情形) 也说不出的确切的东西\footnote{现在知道在这个情形障碍总是为零 [2003 年添加 (MR)：参见 F. Oort，Finite group schemes，local moduli for abelian varieties and lifting problems，Algebraic Geometry Oslo 1970，Wolters-Noordhoff，1972，p. 223--254].}. 希望有人会潜心研究这个问题，它似乎特别有趣. 它与代数结构的``模理论'' (moduli theory) 尤其密切相关.

\begin{corollary}\label{cor:III.6.7}
假设 $H^2(X_0, \mathcal{G}) = 0$，则 $X_{n+1}$ 存在，且若还有 $H^1(X_0, \mathcal{G}) = 0$ 则在同构意义下唯一.
\end{corollary}

\begin{proof}
特别地，逐步进行 (并注意仿射概形对拟凝聚层是无轮回 (acyclic) 的)，得出结论：
\end{proof}

\begin{corollary}\label{cor:III.6.8}
在定理 \ref{thm:III.4.1} 的条件下，若 $X_0$ 仿射，则存在按 $X_0$ 约化的、在 $Y$ 上光滑的 $X$，且这个 $X$ 在同构 (不唯一) 意义下唯一.
\end{corollary}

注意定理 \ref{thm:III.4.1} 的直接证明不能给我们这个结果.

\begin{corollary}\label{cor:III.6.9}
在 \ref{thm:III.6.3} 的条件下，假设 $S$ 仿射，环为 $A$，$\mathcal{I}$ 由 $A$ 的理想 $I$ 定义，最后诸 $\mathrm{gr}_{\mathcal{I}}^{n}(\mathcal{O}_S) = \mathcal{I}^n/\mathcal{I}^{n+1}$ 局部自由. 则 $H^i(X_0, \mathcal{G})$ 等同于 $H^i(X_0, \mathcal{G}_0) \otimes_A \mathrm{gr}_{I}^{n+1}(A)$，其中
\[
\mathcal{G}_0 = \mathfrak{g}_{X_0/S_0},
\]
故 $X_n$ 的延拓障碍类位于 $H^2(X_0, \mathcal{G}_0) \otimes_A \mathrm{gr}_{I}^{n+1}(A)$，且若它为零，解的 (同构) 类集合是 $H^1(X_0, \mathcal{G}_0) \otimes_A \mathrm{gr}_{I}^{n+1}(A)$ 下的主齐性空间.
\end{corollary}

\begin{proof}
特别地：
\end{proof}

\begin{corollary}\label{cor:III.6.10}
在 \ref{cor:III.6.9} 的条件下假设
\[
H^2(X_0, \mathfrak{g}_{X_0/S_0}) = 0,
\]
则在 $S$ 的形式 $\mathcal{I}$-adic 完备化 $\widehat{S}$ 上存在 $\mathcal{I}$-adic 形式概形 $X$，它``在 $S$ 上光滑''(i.e. 诸 $\mathfrak{X}_p$ 在诸 $S_p$ 上光滑) 且按 $X_n$ 约化，i.e. 配备同构
\[
\mathfrak{X} \times_S S_n \stackrel{\sim}{\longleftarrow} X_n.
\]
若还有 $H^1(X_0, \mathfrak{g}_{X_0/S_0}) = 0$，则这样的 $X$ 在同构意义下唯一.
\end{corollary}

\begin{proof}
事实上，逐步构造 $X_{n+1}$，$X_{n+2}$ 等，取诸 $X_i$ 的归纳极限得到 $X$. 唯一性断言已见于前一节.
\end{proof}

\section{在完备局部环 \texorpdfstring{$A$}{A} 上构造光滑形式概形与普通概形的应用}

前一节的结果有时能证明：在一个这样的环上，存在按 $k$ 上给定光滑概形 $X_0$ 约化的 $m$-adic 形式概形. 区分两种情形：

\begin{enumerate}[a)]
\item $A$ 是``等特征''的 (当 $k$ 特征为 0 时尤其如此). 此时知道 $A$ 存在代表元子域 (subfield of representatives)，i.e. 子域 $k'$ 使得 $A \to k$ 诱导同构 $k' \xrightarrow{\sim} k$. 则甚至存在按 $X_0$ 约化的 $A$ 上光滑普通概形，即 $X = X_0 \otimes_k A$，其中 $A$ 借助由 $k'$ 定义的同态 $k \to k' \to A$ 看作 $k$ 上的代数. 但必须注意这个构造不是``自然的''；容易相信 (已经可以在 $A = k[t]/(t^2)$，对偶数代数的情形) 另一个提升同态 $k \to A$ (此处由 $k$ 到自身的绝对导子定义) 定义 $A$ 上的 $X'$，它一般不同构于 $X$ (若 $H^1(X_0, \mathfrak{g}_{X_0/k}) \neq 0$). 研究 (对特征 0 或非完美特征 $p > 0$ 的 $k$) 这样得到的 $A$ 上光滑 $X$ 是哪些、两个同态 $k \to A$ 在什么条件下定义同构的 $A$-概形，将是很有趣的. 然而，$k'$ 的存在足以推出 $X_0$ 提升的第一个障碍 (位于 $H^2\bigl(X_0, \mathfrak{g}_{X_0/k}\bigr) \otimes_k \mathfrak{m}/\mathfrak{m}^2$) 必然为零. 当然，把 $X_0$ 提升为 $A/\mathfrak{m}^2$ 上光滑的 $X_1$ 后，构造 $X_2$ 的新障碍一般不为零：它依赖于 $H^1(X_0, \mathcal{G}_0) \otimes \mathfrak{m}/\mathfrak{m}^2$ 之下某个主齐性空间中一个可变元素，且位于 $H^2(X_0, \mathcal{G}_0) \otimes \mathfrak{m}^2/\mathfrak{m}^3$：应当详细研究这一情形\footnote{它大概由 Kodaira-Spencer 的括号运算描述 (参见 Séminaire Cartan，1960/61，Exp. 4).}.
\item $A$ 是不等特征的. 此时几乎一无所知 (除非碰巧 $H^2(X_0, \mathfrak{g}_{X_0/k}) = 0$，此时可以构造按 $k$ 约化的 $A$ 上 $m$-adic 光滑形式概形). 即使 $A = \mathbf{Z}/p^2\mathbf{Z}$ 且 $X_0$ 是维数 2 的``阿贝尔''概形，也不知道能否把它提升为 $A$ 上光滑的 $X = X_1$\footnote{现在这已被证明，参见第 67 页注 7.}；另一方面，没有例子说明某个 $X_0$ 已被证明不来自 $A$ 上光滑普通概形 $X$. (我觉得这应当存在，取 $X_0$ 为射影曲面)\footnote{这样的例子后来由 J.-P. Serre 构造 (Proc. Nat. Acad. Sc. USA 47 (1961)，No 1，p. 108--109)，至少在若干维数. D. Mumford 给出了一个例子 (未发表)，用代数曲面.}. 只指出，由 Cohen 定理，存在剩余域为 $k$ 的 Cohen $p$-环 $B$ 以及诱导剩余域上恒等同构的同态 $B \to A$；因此，最强的提升结果应取 $A$ 为 Cohen $p$-环：若在一个这样的环之上存在 (普通或形式的) 解，则在剩余域为 $k$ 的每个完备局部环之上也存在. 特别地，由于对 Cohen $p$-环，$\mathfrak{m}/\mathfrak{m}^2$ 典范等同于 $k$，看到对特征 $p > 0$ 的域 $k$ 上的每个光滑概形 $X_0$，存在 $H^2\bigl(X_0, \mathfrak{g}_{X_0/k}\bigr)$ 中的上同调类，即把 $X_0$ 提升为 Cohen $p$-环上光滑概形的第一个障碍；不知道它能否非零\footnote{它可以非零，如第 69 页注 10 所指出的.}.
\end{enumerate}

即使能逐步构造按 $X_0$ 约化的 $X_n$，一般也只能得到按 $X_0$ 约化的、$A$ 上光滑的形式概形 $X$. 当 $X_0$ 在 $A$ 上紧合时，还剩下问题：$X$ 是否事实上可代数化，以便得到按 $X_0$ 约化的、$A$ 上紧合且简单的普通概形. 唯一已知的判则 (在 Bourbaki 讨论班中提及，见于《原理》第三章 4.7.1) 如下：若 $X$ 在 $A$ 上紧合，且 $\mathcal{L}$ 是 $X$ 上的可逆层，使得在 $X_0$ 上诱导的层 $\mathcal{L}_0$ 丰沛 (ample) (i.e. 适当的张量幂 $\mathcal{L}_0^{\otimes n}$ ($n > 0$) 来自 $X_0$ 的射影浸入)，则存在 $A$ 上射影的概形 $X$ 和 $X$ 上的丰沛可逆层 $\mathcal{L}$，使得 $(X, \mathcal{L})$ 经 $m$-adic 完备化得到 $(\mathfrak{X}, \mathcal{L})$. 于是，给定 $X_0$ 上的局部自由层 $\mathcal{E}_0$ (对我们的目的选择丰沛可逆的)，要把它延拓为 $X$ 上的局部自由层 $\mathcal{E}$. 为此归结为在诸 $X_n$ 上逐步构造局部自由层 $\mathcal{E}_n$. 讨论完全类似于 §6 的讨论 (参见注记 \ref{rem:III.6.4})，基本角色由在 $\mathcal{E}_n$ 上诱导恒等的 $\mathcal{E}_{n+1}$ 的自同构层扮演：立即证明该层等同于
\begin{equation}
\mathcal{G} = \operatorname{Hom}_{\mathcal{O}_{X_0}}\bigl(\mathcal{E}_0, \mathcal{E}_0 \otimes \mathrm{gr}_{\mathcal{I}}^{n+1}(\mathcal{O}_X)\bigr) = \operatorname{Hom}_{\mathcal{O}_{X_0}}\bigl(\mathcal{E}_0, \mathcal{E}_0\bigr) \otimes \mathrm{gr}_{\mathcal{I}}^{n+1}(\mathcal{O}_X), \tag{($*$)}
\end{equation}
它仍是交换群层. 得到：

\begin{proposition}\label{prop:III.7.1}
设 $S$ 是配备拟凝聚理想层 $\mathcal{I}$ 的预概形，$X$ 是 $S$ 上的预概形，$S_n$ 是由 $\mathcal{I}^{n+1}$ 定义的 $S$ 的子预概形，$X_n = X \times_S S_n$ (对每个整数 $n$). 设 $\mathcal{E}_n$ 是 $X_n$ 上的局部自由层，要把它延拓为 $X_{n+1}$ 上的局部自由层 $\mathcal{E}_{n+1}$. 则 $\mathcal{E}_n$ 定义 $H^2(X_0, \mathcal{G})$ 中的典范障碍类 (其中 $G$ 是上式给出的拟凝聚层)，该类消没是存在延拓 $\mathcal{E}_n$ 的 $\mathcal{E}_{n+1}$ 的必要且充分条件. 若该类为零，则 (在 $\mathcal{E}_n$ 上诱导恒等的) 同构意义下的解类 $\mathcal{E}_{n+1}$ 的集合，是 $H^1(X_0, \mathcal{G})$ 下的主齐性空间.
\end{proposition}

这个命题给出通常的推论. 只指出，若 $X$ 在 $S$ 上平坦，则可以写
\[
\mathcal{G} = \operatorname{Hom}_{\mathcal{O}_{X_0}}\bigl(\mathcal{E}_0, \mathcal{E}_0\bigr) \otimes_{\mathcal{O}_{S_0}} \mathrm{gr}_{\mathcal{I}}^{n+1}(\mathcal{O}_S),
\]
从而，若 $S$ 仿射环为 $A$ 且诸 $\mathcal{I}^n/\mathcal{I}^{n+1}$ 局部自由，得充分条件
\[
H^2(X_0, \mathcal{G}_0) = 0, \quad \text{其中} \quad \mathcal{G}_0 = \mathbf{Hom}_{\mathcal{O}_{X_0}}\bigl(\mathcal{E}_0, \mathcal{E}_0\bigr),
\]
它保证 $\mathcal{E}_{n+1}$ 存在，从而逐步保证逐次延拓 $\mathcal{E}_m$ ($m = n, n+1$，等) 存在.

回到初始情形，得到：

\begin{proposition}\label{prop:III.7.2}
设 $A$ 是完备局部环，$X$ 是 $A$ 上紧合且平坦的形式概形，使得 $X_0$ 射影且 $H^2(X_0, \mathcal{O}_{X_0}) = 0$. 则存在 $A$ 上射影的概形 $X$，其 $m$-adic 形式完备化同构于 $X$.
\end{proposition}

与 \ref{cor:III.6.10} 结合，得到：

\begin{theorem}\label{thm:III.7.3}
设 $A$ 是剩余域为 $k$ 的完备局部环，$X_0$ 是 $k$ 上射影且光滑的概形，使得
\[
H^2(X_0, \mathfrak{g}_{X_0/k}) = H^2(X_0, \mathcal{O}_{X_0}) = 0.
\]
则存在按 $X_0$ 约化的、$A$ 上光滑且射影的概形 $X$.
\end{theorem}

更一般地，若给出按 $X_0$ 约化的、在 $A_n = A/\mathfrak{m}^{n+1}$ 上光滑的 $X_n$，则存在 $A$ 上光滑且紧合的 $X$ 以及同构 $X \otimes_A A_n = X_n$.

\begin{corollary}\label{cor:III.7.4}
$k$ 上的每条光滑紧合曲线都经约化来自 $A$ 上的光滑紧合曲线.
\end{corollary}

\begin{proof}
正是这个结果 (连同形式几何中层存在性定理) 将是经由超验途径研究 $X_0$ 的基本群的基本工具.
\end{proof}

